\chapter{Background and Literature Review}
\label{ch:background}

This chapter reviews the theoretical foundations of traffic flow prediction (TFP) and explores the evolution of forecasting methodologies from traditional statistical models to advanced deep learning architectures. It subsequently examines the paradigm shift toward Federated Learning (FL) within Intelligent Transportation Systems (ITS), focusing on the two challenges most relevant to this thesis: data heterogeneity caused by non-IID traffic distributions, and communication overhead in bandwidth-constrained edge networks. The review draws upon seminal works and recent advancements to contextualize the contributions of this thesis.

\section{Traditional and Deep Learning Approaches to Traffic Flow Prediction}
Traffic flow prediction has been a cornerstone of intelligent transportation planning for decades. Early research predominantly relied on centralized machine learning and statistical methods, such as combining Principal Component Analysis (PCA) with Support Vector Regression (SVR) \cite{jin2008}, utilizing fuzzy C-mean neural networks \cite{tang2017}, and autoregressive models like ARIMA \cite{arima1970} and Vector Auto-Regression (VAR) \cite{var2005}. While these models laid the groundwork for data-driven traffic management, they assumed data stationarity and struggled to capture the complex, non-linear spatiotemporal dependencies inherent in modern urban traffic networks.

With the proliferation of artificial intelligence, deep learning became the dominant approach. Spatial features from grid-based city structures were extracted using Convolutional Neural Networks (CNN) \cite{zhang2017}, while sequence models like Long Short-Term Memory (LSTM) \cite{hochreiter1997} and Gated Recurrent Units (GRU) \cite{cho2014} captured long-term temporal dependencies in time-series traffic data. Notably, GRU-based Encoder-Decoder (Seq2Seq) architectures have proven effective for multi-step traffic forecasting, as they encode historical observations into a fixed context vector and decode predictions autoregressively across the forecast horizon \cite{zhu2025cmba}.

To model the spatial structure of road networks more explicitly, researchers introduced Spatio-Temporal Graph Neural Networks (STGNN). Li et al. \cite{li2018} proposed the Diffusion Convolutional Recurrent Neural Network (DCRNN), which modeled traffic flow as a diffusion process on a directed graph and captured spatial correlations via bidirectional random walks. Guo et al. \cite{guo2019} introduced ASTGCN, which applied attention mechanisms to capture dynamic spatial and temporal correlations simultaneously. Qi et al. \cite{qi2023fedagcn} proposed the Asynchronous Graph Convolutional Network (AGCN), which extended standard graph convolution by considering asynchronous spatial--temporal dependencies where the influence of traffic conditions at neighboring nodes propagates with a time delay rather than instantaneously. These STGNN-based methods achieve state-of-the-art accuracy in centralized settings but require all sensor data to be consolidated at a single server, an assumption that is increasingly impractical for large-scale urban deployments.

\section{Federated Learning in Intelligent Transportation Systems}
Despite the success of centralized deep learning, such architectures require continuous uploading of raw sensor data to a cloud server, which underutilizes edge computing resources and introduces significant communication overhead, particularly as the number of sensors and model complexity scale up. Federated Learning (FL), introduced by McMahan et al. \cite{mcmahan2017}, addresses these limitations by allowing edge devices to collaboratively train a shared model without exchanging raw data. Instead, only model parameters or gradients are communicated to a central aggregation server.

Recent surveys by Chong et al. \cite{chong2025} and Kairouz et al. \cite{kairouz2021} highlight FL's critical role in ITS applications while identifying two core challenges that remain open: (1) \textit{statistical heterogeneity}, where the non-IID nature of traffic data causes model divergence; and (2) \textit{communication overhead}, where frequent transmission of large model updates strains bandwidth-constrained vehicular networks. The following subsections review the literature addressing each challenge.

\subsection{Data Heterogeneity and Non-IID Traffic Distributions}
A core challenge in FL is statistical heterogeneity, where data across clients is non-Independent and Identically Distributed (non-IID). In the context of traffic prediction, this heterogeneity is especially pronounced: highway sensors observe high-speed, low-variance flows, while urban intersection sensors exhibit stop-and-go patterns with frequent congestion peaks. When a single global model trained via standard FedAvg attempts to fit all these divergent distributions, it often suffers from \textit{client drift}---where local updates pull the global model in conflicting directions---leading to slow convergence and suboptimal accuracy.

Several general-purpose FL algorithms have been proposed to mitigate client drift. Li et al. proposed FedProx \cite{li2020fedprox}, which adds a proximal regularization term to the local objective function, penalizing excessive deviation from the global model. Karimireddy et al. introduced SCAFFOLD \cite{karimireddy2020}, which utilizes control variates to estimate and correct for drift at each client. Li et al. \cite{li2021moon} proposed MOON, which applied model-contrastive learning to align local representations with the global model. Cha and Chang \cite{cha2025} further advanced this line by dynamically adjusting regularization strengths based on MAP-estimated data distributions, achieving significant accuracy improvements in severely skewed non-IID scenarios.

In the ITS domain specifically, Dai et al. \cite{dai2024} proposed a personalized FL framework based on client reputation, where model aggregation weights are assigned proportional to each client's training quality (measured by validation loss), ensuring that clients with more representative or higher-quality data contribute more to the shared model. Li et al. \cite{li2025fedgdan} proposed FedGDAN, a Federated Graph Diffusion Attention Network that introduced an Adaptive Local Aggregation (ALA) mechanism specifically to address non-IID traffic distributions across different road segments. Zhang et al. \cite{zhang2024fedgl} proposed a position-aware structural knowledge sharing framework for federated graph learning in ITS, which decouples structural and attribute features to enable cross-domain knowledge transfer despite heterogeneous data distributions.

\subsection{Clustering-Based Approaches to Data Heterogeneity}\label{sec:clustering}
Rather than regularizing away heterogeneity, an alternative line of research embraces it by grouping clients with similar data distributions into \textit{clusters}, allowing each cluster to train a specialized model. This approach is particularly well-suited for traffic prediction, where sensors naturally group into categories with homogeneous temporal behavior (e.g., highways, arterial roads, residential streets).

Feng et al. \cite{feng2024} proposed FedTFP, a two-stage federated traffic flow prediction model that first partitions the road network into sub-networks using Dynamic Time Warping (DTW) combined with K-means clustering. By ensuring that sensors within the same cluster share similar traffic flow patterns, their approach achieved measurably better performance than random or topology-based partitioning strategies (such as Louvain community detection). Their ablation study demonstrated that pattern-similarity-based partitioning (PSS) was a critical contributor to prediction quality, validating DTW as an effective similarity metric for traffic time series.

Li et al. \cite{li2025fedstdn} adopted a related clustering strategy in their FedSTDN framework for wireless traffic prediction, using data augmentation-based clustering to group cellular base stations into regions with similar traffic patterns before applying federated aggregation within each cluster. This approach reduced communication overhead while achieving a $\sim$32.8\% MSE improvement over non-clustered baselines. Qi et al. \cite{qi2023fedagcn} proposed FedAGCN, which partitioned the traffic network into $K$ sub-graphs and applied federated learning across the sub-graph models, selectively excluding graph convolution parameters from global aggregation to prevent spatial structure conflicts.

Yaqub et al. \cite{yaqub2025flagcn} proposed FLAGCN, which combined asynchronous graph convolutional networks with federated learning by strategically partitioning the traffic network into sub-graphs guided by spatial distance constraints, optimizing the balance between sub-graph count and parameter efficiency to achieve high prediction accuracy with reduced training and inference time.

These works collectively establish that \textit{clustering by traffic pattern similarity is a more principled basis for federated grouping than random or geographic partitioning}, a finding that directly motivates the DTW-based hierarchical approach adopted in this thesis.

\subsection{Communication-Efficient Federated Aggregation}
Transmitting large deep learning models, such as Sequence-to-Sequence (Seq2Seq) GRUs, over bandwidth-constrained vehicular networks remains a significant bottleneck for FL deployments. The literature addresses this challenge through three complementary strategies: gradient compression, intelligent client selection, and hierarchical aggregation.

\textbf{Gradient Compression.} Gao et al. \cite{gao2024fedce} proposed FedCE, which achieves communication efficiency through Top-K gradient sparsification, transmitting only the most significant gradient components while applying an adaptive aggregation strategy that weights client contributions differently. Their results showed that gradient compression can achieve competitive accuracy with substantially reduced uplink bandwidth on wireless traffic prediction benchmarks.

\textbf{Client Selection.} Rather than requiring all clients to participate in every round, intelligent selection strategies reduce communication volume by choosing a representative subset. The Oort framework \cite{lai2021oort} prioritizes clients based on data utility and system speed. Albelaihi \cite{albelaihi2025} proposed MACS, a mobility-aware client selection strategy that forecasts client connectivity and resource availability to optimize participation under latency constraints in IoT environments. In the FL-for-shipping domain, Chen et al. \cite{chen2025bafls} designed a Top-K active learning client selection strategy that identifies the most informative clients each round, improving convergence quality while reducing participation rates.

\textbf{Hierarchical and Multi-Server Aggregation.} Decentralized architectures reduce the single-server bottleneck by distributing aggregation across multiple levels. Ma et al. \cite{ma2025asafl} proposed AS-AFL, a hybrid framework that performs intra-group peer-to-peer aggregation and inter-group asynchronous vertical federated learning, achieving robust performance across over 1000 nodes with no single point of failure. Mazloomi et al. \cite{mazloomi2025} introduced a multi-server FL framework where edge-based local servers perform initial aggregation before inter-server knowledge exchange, reducing uplink latency for mobile vehicular clients. Wang et al. \cite{wang2024cedfl} proposed CEDFL, a communication-efficient decentralized FL mechanism where workers communicate with neighbors according to adaptively determined probabilities, reducing bandwidth consumption by approximately 55\% compared to deterministic topology baselines.

In the specific domain of traffic flow prediction, Zhu et al. \cite{zhu2025cmba} introduced CMBA-FL, which paired a GRU-based Encoder-Decoder architecture with a parameter relevance detector that filters out irrelevant model parameters before upload, achieving a strong accuracy--communication tradeoff on the METR-LA dataset.

\subsection{Federated Learning for Traffic Flow Prediction}
Building on the general FL techniques above, several works have specifically targeted traffic flow prediction in federated settings. Liu et al. \cite{liu2020fedgru} proposed FedGRU, one of the earliest FL-based traffic prediction frameworks, which achieved over 90\% prediction accuracy while preserving data locality. Meng et al. \cite{meng2021cnfgnn} introduced CNFGNN, a cross-node federated graph neural network that enables spatial information sharing across nodes without centralizing data.

More recent work has pushed the boundary further. Wang et al. \cite{wang2024fedgm} proposed FedGM, a Graph Federated Meta-learning strategy that combines meta-learning with federated learning to address topological heterogeneity across different traffic networks, enabling rapid adaptation to unseen graph structures with minimal communication rounds.

\section{Research Gap and Thesis Positioning}
\label{sec:gap}
The literature review reveals that while significant progress has been made in both federated learning algorithms and traffic flow prediction architectures, a key gap persists: \textbf{most FL-based traffic prediction frameworks apply uniform aggregation across all sensors}, ignoring the inherent grouping structure of traffic patterns. Standard FedAvg treats a highway sensor's model update with the same weight as a residential street sensor's, leading to a global model that fits no distribution well. While clustering-based approaches like FedTFP \cite{feng2024} and FedAGCN \cite{qi2023fedagcn} have demonstrated the value of grouping, they operate with \textit{isolated} clusters that do not share knowledge between groups, limiting generalization.

Building upon the GRU-based Encoder-Decoder architecture validated by CMBA-FL \cite{zhu2025cmba}, this thesis proposes a \textit{hierarchical} approach that addresses this gap on two levels:

\begin{enumerate}
    \item \textbf{Intra-cluster specialization:} Sensors are grouped by temporal pattern similarity using DTW-based clustering \cite{feng2024}, and each cluster trains a specialized model via FedAvg with quality-weighted aggregation \cite{dai2024}.
    \item \textbf{Inter-cluster generalization:} Cluster models periodically exchange knowledge through a distance-weighted hierarchical aggregation step, enabling beneficial cross-pattern transfer without collapsing model diversity.
\end{enumerate}

\noindent Combined with adaptive client selection to reduce per-round communication, this framework directly addresses the data heterogeneity and communication challenges identified above, providing a scalable and efficient solution for federated traffic forecasting.
